% s-148-p148-m005; original PDF 148, 149
\Aphorism{s-148-p148-m005}{237}
Im Grunde ist es das fatale Sichaussprechen und Sichausschreiben, das unsere Krankheiten verewigt; die einfachste Vorsicht müsste uns zu gewissen Zeiten Mund und Feder verbieten. Man erleichtert sich für den Augenblick und häuft Infectionsstoff an, an dem man später vielleicht wieder erkrankt: das ist eine gar nicht so unwahrscheinliche Geschichte. Beim Schreiben ist noch eine besondere Gefahr, die dem gesprochenen Worte weniger anhaftet. Der Schreibende, gesetzt, dass er Artist sei, bleibt selten beim einfachen Geständniss, sondern stilisirt, amplificirt, accentuirt; er hat immer neue, immer stärkere Wendungen für sein Erlebniss, immer gewaltsamere Gebärden für sein Pathos. Damit treibt er sein Gefühl selber in die Höhe; seine Ausdrucksweise ist von nun an das Bestimmende, sie hebt das Empfinden zu sich hinauf — was heute noch poetischer Schmuck war, ist morgen schon so gut wie erlebt. Es ist ein unwiderstehlicher Zauber in Schwarz und Weiss, eine Selbsthypnose, auf die noch lange nicht genug geachtet wird.
\par

Schon das bloss negative Idealisiren, das Weglassen der Nebenzüge und Zufälligkeiten, macht aus den Erlebnissen meist etwas zu Heroisches, das später einmal seine Verführung haben wird; der spätere Mensch hat sich vor dem früheren immer zu schämen — aber mit Unrecht, weil der frühere nur als Pose und Fälschung überliefert ist.
\par

